% \clearpage
\section{Related Work}
\paragraph{Large Vision-Language Model}
Motivated by the success of the pretraining technique in LLMs and VFMs, the multimodal learning research community has recently shifted the research attention to LVLMs~\cite{DBLP:journals/corr/abs-2308-01390,DBLP:journals/corr/abs-2305-03726}. 
% The mainstream advanced LVM mainly consists of three components: a text encoder, an image encoder, and a cross-modal alignment module. In particular, a text encoder is usually a language model (for example, LLAVA and Vicuna13B), and an image encoder is always a vision foundation model (such as VIT). The cross-modal alignment module is used to map the visual content into textual space and equip the text encoder with the ability to understand the visual semantics. 
Contemporary advanced LVLMs predominantly feature three core components: a text encoder, an image encoder, and a cross-modal alignment module~\cite{DBLP:conf/emnlp/RohrbachHBDS18}. Specifically, the text encoder often takes the form of a language model, as seen in examples like LLaMA~\cite{llama} and Vicuna~\cite{vicuna}. Conversely, the image encoder is typically derived from VFMs, such as ViT~\cite{vit}. The function of the cross-modal alignment module is to bridge visual content with textual representation, enhancing the text encoder's capacity to interpret visual semantics. 
% Compared with the text encoder and image encoder, the cross-modal alignment module usually has fewer layers and simpler units, such as linear networks. 
To accomplish visual understanding, LVLMs typically undergo multiple training phases~\cite{MultiModal-GPT,minigpt4,llava15,llava,mPLUG-Owl,InstructBLIP}. 
For instance, \citet{llava} first aligns the image features with the word embeddings of a pre-trained LLM during an initial pre-training stage, and subsequently fine-tunes the LVLM using specialized language-image instruction-following datasets. 
For efficiency enhancement, LVLMs often freeze parameters of the LLM or VFM and are trained with efficient fine-tuned techniques~\cite{mPLUG-Owl,InstructBLIP}, such as LoRA~\cite{lora}.

However, in spite of the considerable advancements made by LVLMs, they consistently grapple with hallucination issues. These issues markedly impact their efficacy across a range of vision-language tasks~\cite{DBLP:conf/emnlp/RohrbachHBDS18}. 
% Consequently, devising methods to assess hallucination in VLMs remains a significant challenge.

% To achieve this, VLMs usually are trained with several steps. For example, \citet{llava} align the image features with the pre-trained LLM word embedding by the first pertaining stage and then finetune the VLM with unique language-image instruction-following data. In addition, to improve efficiency, VLMs usually freeze the parameters of the LLM or Vision foundation model and is equipped with efficient fine-tuning methods, such LORA~\cite{}.

% Despite the significant progress achieved by VLMs, they always suffer from hallucinations, which severely influence their performance on various vision-language tasks. Therefore, how to evaluate hallucination in VLMs push a big challenge.

\paragraph{Vision-language Model Hallucinations and Evaluations}
Though hallucination phenomenons and mitigation methods have been extensively studied in the text generation literature~\cite{ji2023survey, factscore}, it is much less investigated in vision-language models~\cite{InstructBLIP, llava,fgaif,mmss}.
% 
Although there are a few existing works tackling this issue, they mainly focus on the constraint problem setting such as image captioning~\cite{densecap}.
For example, 
\citet{DBLP:conf/emnlp/RohrbachHBDS18} propose caption hallucination assessment with image relevance (CHAIR), which is a popular metric for evaluating object hallucination in sentence-level captions. They also show that popular metrics like METEOR~\cite{meteor} and CIDEr~\cite{vedantam2015cider} do not capture this. 
\citet{Li2023EvaluatingOH} extends CHAIR and proposes ``POPE'', a polling-based query technique for probing objects. Besides, \citet{lovenia2023negative} devised Negative Object Presence Evaluation (NOPE) to quantitatively assess object hallucination through VQA, based on ``POPE''. 
\citet{gunjal2023detecting} further proposed to detect hallucinations in more detailed image captions and investigated utilizing a reward model for mitigating them. \citet{DBLP:journals/corr/abs-2309-04041} introduced an evaluation benchmark that contains more diverse types of questions, such as Yes-or-No and Fill-in-the-Blank. 

Different from all the above, we are the first to propose a \textcolor[rgb]{0,0,0}{reference-free} metric for \textcolor[rgb]{0,0,0}{fine-grained} evaluating the answers in the \textcolor[rgb]{0,0,0}{open-ended visual question-answering setting}, where answers are of free form and can be lengthy.


